\section{Achievement of objectives and requirements}\label{sec:achievement_of_objectives_and_requirements}

This section of the report presents the evaluation of the extent to which the proposed system meets the initial objectives and requirements as defined in Chapter~\ref{ch:problem_specification_and_requirements}.

\subsection{Achievement of objectives}\label{subsec:achievement_of_objectives}

\subsubsection{Web crawler}
The primary objective of the web crawler was to be able to efficiently crawl multiple e-commerce platforms and collect page data.
The designed distributed web crawler was able to crawl over 116,000 pages crawled across 36 different e-commerce websites.
At its current implementation of one instance, each crawling 100 pages, being spun up every 3 minutes, it was able to crawl over 29,000 pages in a single day.

\subsubsection{Boilerplate removal and HTML minimization}
This was a secondary objective of the system, which was to be able to remove boilerplate content and minimize the HTML content of the crawled pages.
The HTML minimizer algorithm was able to reduce the size of the HTML content by \texttt{98.16\%}, and was able to minimize batches of 20,000 documents in an average of \texttt{768.21 seconds}.

\subsubsection{Page classification}
The primary objective was to evaluate different page classification techniques and select the best performing model for deployment.
The techniques tested were:
\begin{itemize}
    \item TF-IDF with traditional machine learning classifiers
    \item BERT embeddings with a neural network classifier
    \item MarkupLM for document classification
\end{itemize}

While the best performing model was the MarkupLM model, the selected model for deployment was the BERT embeddings with a neural network classifier, due to its balance between performance, speed, and cost.
The final trained model for this task achieved an F1-score of \texttt{0.971}, and could classify batches of 20,000 documents in an average of \texttt{40.89 minutes}.

\subsubsection{Information extraction}
The primary objective was to evaluate different information extraction techniques and select the best performing model for deployment.
The techniques tested were:
\begin{itemize}
    \item MarkupLM for token classification
    \item Multi-modal approach with LayoutLMv3
    \item LLM based approach
\end{itemize}

While the LLM based approach showed the best performance, the selected model for deployment was the MarkupLM model, due to its balance between performance, speed, and cost.
The final trained model for this task achieved an overall F1-score of \texttt{0.970}, and could extract information from batches of 20,000 documents in an average of \texttt{28.5 minutes}.

The results from this module also led into the price history tracking framework, which was able to keep records of price changes for a given product at a specific e-commerce platform.

\subsubsection{Product matching}
The primary objective was to evaluate different product matching techniques and select the best performing model for deployment.
10 different techniques were tested, including those using string similarity algorithms such as trigram Jaccard similarity, edit distance, algorithms such as the BM-25 algorithm, and deep learning approaches using embeddings models such as BERT, RexBERT, and MarkupLM.

The best performing technique was to use trigram Jaccard similarity to generate candidate matches, and then use fine-tuned RexBERT embedding model to generate vectors for the minimized page content, and then use cosine similarity to find the best matches.
This technique was able to achieve an F1-score of \texttt{0.88}.

\subsubsection{Proof-of-concept system}
This was a secondary objective of the system, which was to implement a web-based application that allows users to search for products and view their historical prices across different platforms, as well as compare prices and identify the best deals available.
The implemented system was able to retrieve price history for a given product at a specific retailer with a p95 response time of \texttt{224ms}, and was able to retrieve cross-site product matches for a given product with a p95 response time of \texttt{2,907ms}.

\subsection{Achievement of requirements}\label{subsec:achievement_of_requirements}

Table~\ref{tab:requirements-achievement} summarizes the extent to which each functional and non-functional requirement was achieved.

\begin{table}[H]
    \centering
    \small
    \renewcommand{\arraystretch}{1.3}
    \begin{tabularx}{\textwidth}{p{0.48\textwidth} p{0.48\textwidth}}
        \toprule
        \textbf{Requirement} & \textbf{Achievement} \\
        \midrule
        Provided a set of seed URLs, the system must crawl each URL, render the page, and store its HTML content in a BLOB storage. &
        \textbf{Achieved.} Pages are crawled at a rate of \texttt{29,000} pages per day and content is stored in an AWS S3 bucket. \\
        \midrule
        For each page crawled, the system must identify hyperlinks to other pages and add them to the crawling queue. &
        \textbf{Achieved.} The crawler can identify outgoing hyperlinks from a page and add them to the queue. \\
        \midrule
        The system must be able to classify each crawled webpage as a product page or a non-product page. &
        \textbf{Achieved.} The page classification module classifies \texttt{20,000} pages in \texttt{40.89} minutes with an F1 score of \texttt{0.971}. \\
        \midrule
        For each product page, the system must extract the product's title, current price, product image, and whether the product is in stock or out of stock. &
        \textbf{Achieved.} The information extraction module processes \texttt{20,000} pages in \texttt{28.5} minutes with an F1 score of \texttt{0.970}. \\
        \midrule
        The system must be able to display the price-history of a product over time. &
        \textbf{Achieved.} Using CDC on the extraction information table, the system tracks changing prices. \\
        \midrule
        Given a specific product from a specific e-commerce website, the system must be able to find the same product offered on other e-commerce websites. &
        \textbf{Achieved.} The product matching module achieves an F1 score of \texttt{0.88}. \\
        \midrule
        The system must be able to track at least 50,000 products from at least 30 different e-commerce websites. &
        \textbf{Achieved.} The system tracks \texttt{73,354} products across \texttt{36} e-commerce websites. \\
        \midrule
        The price of each product must be updated at least once every 72 hours. &
        \textbf{Achieved.} \texttt{97.4\%} of all products were updated within \texttt{72} hours. \\
        \midrule
        The distributed crawler must ensure that the crawling process does not impact the performance of the target e-commerce websites, and must respect the robots.txt file of each website. &
        \textbf{Partially achieved.} This is difficult to measure directly because infrastructure metrics are not available for target websites.
        However, backoff mechanisms are implemented to minimize impact on target servers. \\
        \midrule
        The system must be able to retrieve the price history for a product in under 500 milliseconds. &
        \textbf{Achieved.} p95 response time for price history retrieval was \texttt{224ms}. \\
        \midrule
        The system must be able to find the same product offered on other e-commerce websites in under 2 seconds. &
        \textbf{Not achieved.} p95 response time for similar product retrieval was \texttt{2,907ms}. \\
        \midrule
        The system must be implemented at a low cost of under USD 300 per month, including cloud hosting and storage costs. &
         \\
        \bottomrule
    \end{tabularx}
    \caption{Achievement of project functional and non-functional requirements}
    \label{tab:requirements-achievement}
\end{table}